\begin{figure}[htbp]
    \centering
    \begin{tikzpicture}[
        every node/.style = {
            draw, 
            rectangle, 
            rounded corners=3pt, 
            align=center, 
            minimum height=0.75cm, 
            minimum width=2.2cm, 
            font=\sffamily\footnotesize,
            inner sep=3pt
        },
        root/.style = {fill=blue!12, draw=blue!80!black, thick},
        latent/.style = {fill=orange!15, draw=orange!80!black, thick},
        evidence/.style = {fill=green!15, draw=green!60!black, thick, dashed},
        target/.style = {fill=red!15, draw=red!80!black, thick},
        arrow/.style = {->, >=stealth, thick, draw=gray!80}
    ]

    % Level 0: Astrophysical Priors (Root Nodes)
    \node[root] (X0) at (-2.4, 0.0) {$X_0$\\Stellar M-Dwarf};
    \node[root] (X1) at ( 2.4, 0.0) {$X_1$\\Orbital HZ};
    
    % Level 1: Primary Latent State & Observable H2O
    \node[latent]   (X2) at ( 0.0, -1.8) {$X_2$\\Hycean Planet};
    \node[evidence] (X6) at ( 3.7, -1.8) {$X_6$\\Spectro H$_2$O};
    
    % Level 2: Secondary Latent State & Atmospheric Biomarkers
    \node[evidence] (X4) at (-3.7, -3.6) {$X_4$\\Spectro CH$_4$};
    \node[latent]   (X3) at ( 0.0, -3.6) {$X_3$\\Biological Ocean};
    \node[evidence] (X5) at ( 3.7, -3.6) {$X_5$\\Spectro CO$_2$};
    
    % Level 3: Target Astrobiological Hypotheses
    \node[target]   (X7) at (-1.8, -5.3) {$X_7$\\Biomarker DMS};
    \node[target]   (X8) at ( 1.8, -5.3) {$X_8$\\Techno CFC};

    % Causal Directed Edges
    \draw[arrow] (X0) -- (X2);
    \draw[arrow] (X1) -- (X2);
    \draw[arrow] (X1) -- (X6);
    
    \draw[arrow] (X2) -- (X3);
    \draw[arrow] (X2) -- (X4);
    \draw[arrow] (X2) -- (X5);
    
    \draw[arrow] (X3) -- (X7);
    \draw[arrow] (X3) -- (X8);

    % Symmetrical, Centered Horizontal Legend underneath
    \begin{scope}[shift={(0.0, -6.6)}]
        \node[root, minimum width=0.45cm, minimum height=0.35cm] (l1) at (-4.2, 0) {};
        \node[draw=none, right=0.08cm of l1, font=\sffamily\scriptsize, anchor=west] {Prior / Root};
        
        \node[latent, minimum width=0.45cm, minimum height=0.35cm] (l2) at (-1.5, 0) {};
        \node[draw=none, right=0.08cm of l2, font=\sffamily\scriptsize, anchor=west] {Latent State};
        
        \node[evidence, minimum width=0.45cm, minimum height=0.35cm] (l3) at ( 1.1, 0) {};
        \node[draw=none, right=0.08cm of l3, font=\sffamily\scriptsize, anchor=west] {Evidence (JWST)};
        
        \node[target, minimum width=0.45cm, minimum height=0.35cm] (l4) at ( 4.0, 0) {};
        \node[draw=none, right=0.08cm of l4, font=\sffamily\scriptsize, anchor=west] {Target Hypothesis};
    \end{scope}

    \end{tikzpicture}
    \caption{\textbf{Causal Topology of the K2-18b Astrobiological Bayesian Network.} Directed Acyclic Graph (DAG) representing the 9-node formulation used in this thesis. Nodes are structured hierarchically: root astrophysical priors ($X_0, X_1$), latent planetary states ($X_2, X_3$), spectroscopic evidence observable by JWST ($X_4, X_5, X_6$), and target biological/technological hypotheses ($X_7, X_8$).}
    \label{fig:k218b_dag}
\end{figure}
